% Standalone problem document, generated from the adjacent README.md.
% Compile with XeLaTeX. Mathematical content is maintained in Markdown.
\documentclass[11pt,a4paper]{article}
\usepackage[fontsize=10.5pt]{scrextend}
\usepackage[margin=22mm,headheight=15pt,headsep=9mm,footskip=11mm]{geometry}
\usepackage{fontspec}
\setmainfont{texgyrepagella}[Extension=.otf,UprightFont=*-regular,BoldFont=*-bold,ItalicFont=*-italic,BoldItalicFont=*-bolditalic]
\setsansfont{texgyreheros}[Extension=.otf,UprightFont=*-regular,BoldFont=*-bold,ItalicFont=*-italic,BoldItalicFont=*-bolditalic]
\setmonofont{texgyrecursor}[Extension=.otf,UprightFont=*-regular,BoldFont=*-bold,ItalicFont=*-italic,BoldItalicFont=*-bolditalic,Scale=MatchLowercase]
\usepackage{amsmath,amssymb}
\usepackage{unicode-math}
\setmathfont{texgyrepagella-math.otf}
\usepackage{xcolor,microtype,enumitem,needspace,fancyhdr}
\usepackage{bookmark}
\definecolor{ink}{HTML}{17344B}
\definecolor{accent}{HTML}{197D80}
\definecolor{muted}{HTML}{596773}
\hypersetup{colorlinks=true,linkcolor=accent,urlcolor=accent,pdftitle={SF-01 - Preservation
of real H-matrix structure by Newton square-root
iterates},pdfauthor={Open Problems in Numerical Linear Algebra}}
\urlstyle{same}
\setlength{\parindent}{0pt}
\setlength{\parskip}{5pt plus 1pt minus 1pt}
\linespread{1.04}
\setlength{\emergencystretch}{3em}
\widowpenalty=10000
\clubpenalty=10000
\displaywidowpenalty=10000
\allowdisplaybreaks[1]
\setlist{nosep,leftmargin=1.5em}
\providecommand{\tightlist}{\setlength{\itemsep}{0pt}\setlength{\parskip}{0pt}}
\providecommand{\pandocbounded}[1]{#1}
\pagestyle{fancy}
\fancyhf{}
\fancyhead[L]{\sffamily\footnotesize\color{muted}OPEN PROBLEMS IN NUMERICAL LINEAR ALGEBRA}
\fancyhead[R]{\sffamily\bfseries\footnotesize\color{accent}SF-01}
\fancyfoot[L]{\parbox[b]{0.9\textwidth}{\sffamily\fontsize{8}{10}\selectfont\color{muted}Editorial ratings; a bounded search cannot certify that a problem remains unsolved.\\Literature check: 2026-09-08}}
\fancyfoot[R]{\sffamily\footnotesize\color{muted}\thepage}
\renewcommand{\headrulewidth}{0.3pt}
\renewcommand{\headrule}{\hbox to\headwidth{\color{accent}\leaders\hrule height \headrulewidth\hfill}}
\makeatletter
\renewcommand{\section}{\Needspace{5\baselineskip}\@startsection{section}{1}{0pt}{12pt plus 2pt minus 2pt}{4pt}{\sffamily\bfseries\large\color{ink}}}
\renewcommand{\subsection}{\Needspace{4\baselineskip}\@startsection{subsection}{2}{0pt}{9pt plus 1pt minus 1pt}{3pt}{\sffamily\bfseries\normalsize\color{ink}}}
\makeatother
\setcounter{secnumdepth}{0}
\begin{document}
{\sffamily\small\color{muted}Matrix functions and stability\par}
\vspace{3pt}
{\raggedright\hyphenpenalty=10000\exhyphenpenalty=10000\sffamily\bfseries\fontsize{21}{24}\selectfont\color{ink}Preservation
of real H-matrix structure by Newton square-root iterates\par}
\vspace{7pt}
{\sffamily\small\textbf{Difficulty:} challenging\quad\textbf{Importance:} interesting
to specialist\par}
\vspace{4pt}
{\color{accent}\hrule height 0.8pt}
\vspace{6pt}
\textbf{Status:} open in the cited book; no subsequent resolution
located.

\subsection{Problem statement}\label{problem-statement}

For a real square matrix \(Y=(y_{ij})\), define its comparison matrix by
\(\mathcal M(Y)_{ii}=|y_{ii}|\) and \(\mathcal M(Y)_{ij}=-|y_{ij}|\) for
\(i\ne j\). Call \(Y\) a nonsingular H-matrix when
\(\mathcal M(Y)=sI-B\) for some entrywise nonnegative \(B\) and real
\(s>\rho(B)\).

Let \(n\ge1\) and let \(A\in\mathbb R^{n\times n}\) be a nonsingular
H-matrix with \(a_{ii}>0\) for every \(i\). Consider the
exact-arithmetic recurrence

\[
X_0=A,\qquad X_{k+1}=\tfrac12(X_k+X_k^{-1}A),\qquad k\ge0.
\]

Is every \(X_k\) a nonsingular H-matrix with strictly positive diagonal?
The known existence of these iterates and their convergence to the
principal square root do not themselves establish the requested
structure at every step.

\subsection{References and status}\label{references-and-status}

N. J. Higham,
\href{https://epubs.siam.org/doi/10.1137/1.9780898717778.ch6}{\emph{Functions
of Matrices: Theory and Computation}} (SIAM, 2008), §6.3, equation
(6.12), §6.8.3, and Research Problem 6.25, p.~170. The author's
\href{https://nhigham.com/errata-for-functions-of-matrices-theory-and-computation/}{errata}
correct the surrounding p.~161 theorem to real H-matrices; that
restriction is adopted here. L. Lin and Z.-Y. Liu,
\href{https://doi.org/10.1023/A:1012931928589}{\emph{On the Square Root
of an H-matrix with Positive Diagonal Elements}}, \emph{Annals of
Operations Research} 103 (2001), 339--350, concern the square root
itself. Searches on 2026-09-08 for the problem number and
Newton/H-matrix structure preservation found no resolution. The recent
\href{https://arxiv.org/html/2605.21679}{Bini--Iannazzo--Meini--Meng
preprint} (20 May 2026), §§3--4, treats M-matrices, a narrower class.
The latest explicit open-problem evidence found for this assertion
remains the 2008 book.
\end{document}
